

\subsection{Ficha 1: Seguridad en microservicios}

\begin{description}

\item[Referencia (IEEE):]
\fullcite{jayalath2024}

\item[DOI o ISBN verificado:]
\url{https://doi.org/10.1109/ACCESS.2024.3481374}

\item[Resumen:]
Este artículo aborda los problemas de seguridad que se pueden presentar en las arquitecturas de microservicios, particularmente por la comunicación entre servicios y dependencias que existen entre sus componentes. Para analizar esta situación, los autores realizaron una revisión de 62 estudios académicos mediante la cual identificaron 126 vulnerabilidades y las organizaron en una clasificación. De igual manera, evaluaron cuatro aplicaciones de referencia utilizando tres herramientas de análisis de seguridad. Durante las pruebas detectaron 1667 vulnerabilidades. Los hallazgos destacan la importancia de utilizar diferentes herramientas de detección y mantener una supervisión continua para identificar riesgos y fortalecer la protección de los sistemas basados en microservicios.

\item[Dato concreto:]
Los investigadores detectaron 1667 vulnerabilidades durante el análisis de cuatro aplicaciones de microservicios utilizando tres herramientas de escaneo.

Ubicación: sección V, página 155193.

\item[Nota personal:]
Esta investigación me permitió comprender que la arquitectura de microservicios puede presentar diferentes riesgos de seguridad debido a la comunicación y las dependencias entre sus componentes. En TA-IND-01 utilizaré estos conocimientos para analizar los posibles riesgos y las medidas de protección que deben considerarse en un sistema distribuido como WhatsApp.

\end{description}



\subsection{Ficha 2}
\begin{description}
  \item[Referencia (IEEE):]
  \item[DOI o ISBN verificado:]
  \item[Resumen:]
  \item[Dato concreto:]
  \item[Nota personal:]
\end{description}

\subsection{Ficha 3}
\begin{description}
  \item[Referencia (IEEE):]
  \item[DOI o ISBN verificado:]
  \item[Resumen:]
  \item[Dato concreto:]
  \item[Nota personal:]
\end{description}

\subsection{Ficha 4}
\begin{description}
  \item[Referencia (IEEE):]
  \item[DOI o ISBN verificado:]
  \item[Resumen:]
  \item[Dato concreto:]
  \item[Nota personal:]
\end{description}

\subsection{Ficha 5}
\begin{description}
  \item[Referencia (IEEE):]
  \item[DOI o ISBN verificado:]
  \item[Resumen:]
  \item[Dato concreto:]
  \item[Nota personal:]
\end{description}
